
\hypertarget{sec:derivatives}{%
\section{Tangents on Manifolds}\label{sec:derivatives}}

A standard definition of the directional derivative for a function
\(\boldsymbol{v}\colon \mathcal{U}\rightarrow\mathcal{V}\) between vector spaces
\(\mathcal{U}\) and \(\mathcal{V}\) is \citep{gurtin1981introduction}:

\begin{equation}
\begin{aligned}
D\boldsymbol{v}\cdot\boldsymbol{u} &= \lim _{\varepsilon \rightarrow 0} \frac{\boldsymbol{v}(\boldsymbol{x}+\varepsilon \boldsymbol{u})- \boldsymbol{v}(\boldsymbol{x})}{\varepsilon} \\
&= \left.\frac{\mathrm{d}}{\mathrm{d} \varepsilon}\boldsymbol{v}(\boldsymbol{x} + \varepsilon\boldsymbol{u})\right|_{\varepsilon=0}.
\end{aligned}
\label{eq:dlinear}
\end{equation}
% \FCF{write above equation in a single line?}
However, when the domain
\(\mathcal{U}\) and co-domain \(\mathcal{V}\) are nonlinear manifolds,
both the increment \(\boldsymbol{x} + \varepsilon\boldsymbol{u}\) in \(\mathcal{U}\) and the difference of \(\boldsymbol{v}\) in \(\mathcal{V}\)
are no longer defined and 
the definition in \cref{eq:dlinear} breaks down.
The starting point for  %adopted in 
this work is to assume that for any sufficiently smooth
single-parameter curve
\(\gamma\colon\mathbb{R} \rightarrow \mathcal{U}\) through a manifold
\(\mathcal{U}\), one is able to identify \emph{velocity} vectors
\(\dot{\gamma}(\varepsilon)\) at points \(\varepsilon\) that are
\emph{tangent} to the curve \(\gamma(\varepsilon)\).
A curve \(\gamma\) whose tangent or velocity, \(\dot{\gamma}\), at each point
is given by a vector field \(\boldsymbol{\eta}\) %on
in the tangent bundle\footnote{A \emph{tangent bundle} \(T\mathcal{U}\) is the collection of all tangent spaces \(T_x\mathcal{U}\) for the points \(\boldsymbol{x}\) in \(\mathcal{U}\).} 
\(T\mathcal{U}\)
is known as an \emph{integral curve} of \(\boldsymbol{\eta}\). 
% In practice this is often intuitive, and in all cases considered,
% an integral curve and its tangent will be apparent at the outset.
A formal construction of this notion can be found
in standard geometry texts \citep{spivak1979comprehensive}.

\hypertarget{directional-derivatives}{%
\subsubsection{Scalar-Valued Functions}\label{directional-derivatives}}

For \emph{scalar-valued functions} and \emph{functionals}
\(f\colon \mathcal{U}\rightarrow \mathbb{R}\) on a manifold
\(\mathcal{U}\), the co-domain \(\mathbb{R}\) is naturally linear and
only the increment \(\boldsymbol{x}+\varepsilon \boldsymbol{u}\) of Equation~(\ref{eq:dlinear}) in
the domain \(\mathcal{U}\) is undefined.
In this case the directional derivative at a point \(\boldsymbol{x}\) in \(\mathcal{U}\) in the direction
\(\boldsymbol{u}\) is defined in terms of an integral curve \(\gamma\) with the properties:
\begin{equation}
  \gamma(0)=\boldsymbol{x}
  \qquad\text{ and }\qquad
  \dot{\gamma}(0)=\boldsymbol{u}
  \label{eq:integral-curve}
\end{equation} 
where \(\boldsymbol{u}\) is a vector in the tangent space \(T_{\gamma(0)}\mathcal{U}\).
Then the directional derivative \(D f\cdot\boldsymbol{u}\) of \(f\)
in the direction \(\boldsymbol{u}\) is defined by: 
\begin{equation}
\begin{aligned}
%D f\cdot \boldsymbol{u} &\triangleq \lim _{\varepsilon \rightarrow 0} \frac{f\left(\gamma(\varepsilon)\right)-f\left(\gamma(0)\right)}{\varepsilon} \\
%&=\left.\frac{\mathrm{d}}{\mathrm{d} \varepsilon} f \circ \gamma(\varepsilon)\right|_{\varepsilon=0}
D f\cdot \boldsymbol{u} 
  &\triangleq \lim _{\varepsilon \rightarrow 0} \frac{f\bigl(\gamma(\varepsilon)\bigr)-f\bigl(\gamma\left(0\right)\bigr)}{\varepsilon} \\
&=\left.\frac{\mathrm{d}}{\mathrm{d} \varepsilon}  f \circ \gamma(\varepsilon) \right|_{\varepsilon=0}
\end{aligned}
\label{eq:dscalar}
\end{equation}
%
This definition is independent of the
particular choice of \(\gamma\), as long as it satisfies
Equation~(\ref{eq:integral-curve}).
When \(f\) acts on functions,
the manifold \(\mathcal{U}\) is infinite-dimensional and Equation~(\ref{eq:dscalar}) is often
referred to as the \emph{Gateaux} derivative \citep{simo1992symmetric,zienkiewicz2014finite}. However, many authors reserve this term for the more restrictive setting of infinite-dimensional \emph{vector} spaces.

\hypertarget{covariant-derivatives}{%
\subsubsection{Tensor-Valued Functions: Covariant Derivatives}\label{covariant-derivatives}}

Extending \cref{eq:dscalar} to a vector or tensor field
\(\boldsymbol{v}: \mathcal{U}\rightarrow T\mathcal{U}\) on a manifold
\(\mathcal{U}\)
would require
evaluating the limit: 
\begin{equation*}
%D \boldsymbol{v}(\chi)\cdot \boldsymbol{v}=\left.\frac{\mathrm{d}}{\mathrm{d} t}\right|_{t=0} \boldsymbol{v}(\gamma(t))=
%\lim _{\varepsilon \rightarrow 0} \frac{\boldsymbol{v} (\gamma(\varepsilon))-\boldsymbol{v}(\gamma(0))}{\varepsilon} .
\lim _{\varepsilon \rightarrow 0} \frac{\boldsymbol{v} \bigl( \gamma(\varepsilon) \bigr) -\boldsymbol{v}\bigl(\gamma(0)\bigr) }{\varepsilon} .
\end{equation*}
% \FCF{note again the use of brackets for vector function}
This, however, is again ill-defined because
\(\boldsymbol{v} \bigl( \gamma(\varepsilon) \bigr) \in T_{\gamma(\varepsilon)} \mathcal{U}\)
and \(\boldsymbol{v}\bigl( \gamma(0) \bigr) \in T_{\gamma(0)} \mathcal{U}\) belong to
different vector spaces. There are two constructions which circumvent
this: the \emph{Lie derivative} \(\mathcal{L}\), and
\emph{differentiation by a connection} \(\nabla\).

A \emph{connection} \(\nabla\) is an additional structure imparted on a
manifold that explicitly dictates how the tangent space at each point in
the manifold is transported. A connection on a finite-dimensional manifold is identified on a basis
\(\{\mathbf{e}_k\}\) by the Christoffel symbols \(\Gamma^k{ }_{i j}\) for the basis 
\citep{abraham1994foundations, bishop1980tensor}: 
\begin{equation*}
\nabla_{\mathbf{e}_j} \mathbf{e}_i=\Gamma^k{ }_{i j} \mathbf{e}_k
% \label{eq:christoffel}
\end{equation*}
where summation on repeated indices is implied. On a Riemannian
manifold a natural choice for \(\nabla\) is the Levi-Civita connection,
which is defined entirely by the Riemannian metric
\(\left<\cdot,\,\cdot\right>\) and furnishes the \emph{covariant
derivative}. 

\paragraph{Vectors}
Given a point $p$ of the manifold $M$, a vector field $\bm{v}: M \rightarrow T_p M$ defined in a neighborhood of $p$ and a tangent vector $\bm{u} \in T_p M$, the covariant derivative of $\bm{v}$ at $p$ along $\bm{u}$ is the tangent vector at $p$, denoted $\left(\nabla_{\bm{u}} \bm{v}\right)_p$, such that the following properties hold (for any tangent vectors $\bm{u}, \bm{x}$ and $\bm{y}$ at $p$, vector fields $\bm{v}$ and $\bm{w}$ defined in a neighborhood of $p$, scalar values $\alpha$ and $\beta$ at $p$, and scalar function $f$ defined in a neighborhood of $p$ ):
\begin{itemize}
\item $\left(\nabla_{\bm{u}} \bm{v}\right)_p$ is linear in $\bm{u}$
$$
\left(\nabla_{\alpha \bm{x} + \beta \bm{y}} \bm{v}\right)_p
= \alpha\, \left(\nabla_{\bm{x}} \bm{v}\right)_p+\beta\, \left(\nabla_{\bm{y}} \bm{v}\right)_p
$$
\item $\left(\nabla_{\bm{u}} \bm{v}\right)_p$ is additive in $\bm{v}$:
$$
\left(\nabla_{\bm{u}}[\bm{v}+\bm{w}]\right)_p=\left(\nabla_{\bm{u}} \bm{v}\right)_p+\left(\nabla_{\bm{u}} \bm{w}\right)_p
$$
\item $\left(\nabla_{\bm{u}} \bm{v}\right)_p$ obeys the product rule:
$$
\left(\nabla_{\bm{u}}[f \bm{v}]\right)_p=f(p)\left(\nabla_{\bm{u}} \bm{v}\right)_p+\left(\nabla_{\bm{u}} f\right)_p \bm{v}_p
$$
where $\nabla_{\bm{u}} f$ is defined by \cref{eq:dscalar}.
\end{itemize}

This can is expressed in coordinates (Romano):
$$
\nabla_{\bm{u}} \bm{v}=Y^j\left(X_{/ j}^i+\Gamma_{j k}^i X^k\right) \mathbf{e}_i
$$
Proof:
$$
\begin{aligned}
\nabla_{\bm{u}} \bm{v} & =\nabla_{\left(Y^j \mathbf{e}_j\right)}\left(X^k \mathbf{e}_k\right)=Y^j \nabla_{\mathbf{e}_j}\left(X^k \mathbf{e}_k\right)= \\
& =Y^j\left[\left(\partial_{\mathbf{e}_j} X^k\right) \mathbf{e}_k+\left(\nabla_{\mathbf{e}_j} \mathbf{e}_k\right) X^k\right]= \\
& =Y^j\left(X_{/ j}^i+\Gamma_{j k}^i X^k\right) \mathbf{e}_i
\end{aligned}
$$
%
The \emph{directionless} covariant derivative $\nabla \bm{v}$ of a vector field $\bm{v}$ is represented by the second-order tensor field
$$
\nabla \bm{v}=\left(\partial_j v^i+v^k \Gamma_{j k}^{~ i}\right) \bm{\partial}_i \otimes \bm{d} X^j
$$

$$
\begin{aligned}
\nabla \bm{v} & =\nabla_{\bm{\partial}_j} \bm{v} \otimes \bm{d} X^j \\
& =\nabla_{\bm{\partial}_j}\left(v^i \bm{\partial}_i\right) \otimes \boldsymbol{d} X^j \\
%&=\boldsymbol{v}_{; j} \otimes \mathbf{G}^j\\
& =\nabla_{\boldsymbol{\partial}_j}\left(v^i\right) \boldsymbol{\partial}_i \otimes \boldsymbol{d} X^j+v^i \nabla_{\partial_j}\left(\bm{\partial}_i\right) \otimes \boldsymbol{d} X^j \\
& =\partial_j v^i \, \bm{\partial}_i \otimes \boldsymbol{d} X^j+v^i \Gamma_{j i}^{. C} \partial_C \otimes \boldsymbol{d} X^j \\
& =\left(\partial_j v^i+v^k \Gamma_{j k}^{~ i}\right) \boldsymbol{\partial}_i \otimes \boldsymbol{d} X^j,
\end{aligned}
$$
Remarks:
\begin{itemize}
    \item Note that $\left(\nabla_{\mathbf{v}} \mathbf{u}\right)_p$ depends not only on the value of $\mathbf{u}$ at $p$ but also on values of $\mathbf{u}$ in an infinitesimal neighborhood of $p$ because of the last property, the product rule.
\end{itemize}

\paragraph{Co-vectors}
For a \emph{co-vector} field \(\boldsymbol{\omega}\) the covariant
derivative \(\nabla_{\boldsymbol{u}} \boldsymbol{\omega}\) at a point
\(\boldsymbol{x}\in\mathcal{U}\) is defined by: 
\begin{equation}
\left(\nabla_{\boldsymbol{u}} \boldsymbol{\omega}\right) \left[\boldsymbol{\eta}\right]\triangleq D \boldsymbol{\omega}\left[\boldsymbol{\eta}\right] \cdot{\boldsymbol{u}}-\boldsymbol{\omega}\left[\nabla_{\boldsymbol{u}} \boldsymbol{\eta}\right]
\label{eq:dcovector}
\end{equation} 
for any vector field \(\boldsymbol{\eta}\) in
a neighborhood of \(\boldsymbol{x}\). 
Note that the application
\(\boldsymbol{\omega}\left[\boldsymbol{\eta}\right]\) of a covector \(\boldsymbol{\omega}\) to a
vector \(\boldsymbol{\eta}\) is a scalar so that the expression
\(D \boldsymbol{\omega}\left[\boldsymbol{\eta}\right] \cdot{\boldsymbol{u}}\) in
Equation~(\ref{eq:dcovector}) is given by Equation~(\ref{eq:dscalar}).

\subsubsection{Tensor-Valued Functions: Lie Derivatives}

The main difference between the Lie derivative and a derivative with respect to a connection is that the latter derivative of a tensor field with respect to a tangent vector is well-defined even if it is not specified how to extend that tangent vector to a vector field. However a connection requires the choice of an additional geometric structure (e.g. a Riemannian metric or just an abstract connection) on the manifold. 
In contrast, when taking a Lie derivative, no additional structure on the manifold is needed, but it is impossible to talk about the Lie derivative of a tensor field with respect to a single tangent vector, since the value of the Lie derivative of a tensor field with respect to a vector field X at a point \(p\) depends on the value of X in a neighborhood of \(p\), not just at \(p\) itself. 
Finally, the exterior derivative of differential forms does not require any additional choices, but is only a well defined derivative of differential forms (including functions).

\paragraph{Vectors}
$$
[X, Y]_x=\left(\mathcal{L}_X Y\right)_x:=\lim _{t \rightarrow 0} \frac{\left(\mathrm{D} \Phi_{-t}^X\right) Y_{\Phi_t^X(x)}-Y_x}{t}=\left.\frac{\mathrm{d}}{\mathrm{d} t}\right|_{t=0}\left(\mathrm{D} \Phi_{-t}^X\right) Y_{\Phi_t^X(x)}
$$
or
$$
[\bm{X}, \bm{\eta}]_s=\left(\mathscr{L}_{\bm{X}} \bm{\eta}\right)_s\triangleq\left.\frac{d}{d t}\right|_{t=s} \bm{\chi}_{s, t *} \bm{\eta}_t
$$

\paragraph{Tensor}
For every smooth 1-parameter family $\Phi_t$ of diffeomorphisms that integrate a vector field $X$ in the sense that $\left.\frac{d}{d t}\right|_{t=0} \Phi_t=X \circ \Phi_0$, one has
$$
\mathcal{L}_X T=\left.\left(\Phi_0^{-1}\right)^* \frac{d}{d t}\right|_{t=0} \Phi_t^* T=-\left.\frac{d}{d t}\right|_{t=0}\left(\Phi_t^{-1}\right)^* \Phi_0^* T
$$

\paragraph{Forms}
The Cartan formula can be used as a definition of the Lie derivative of a differential form:
$$
\mathcal{L}_X \omega=d\left(\iota_X \omega\right)+\iota_X d \omega .
$$
Cartan's formula shows in particular that
$$
d \mathcal{L}_X \omega=\mathcal{L}_X(d \omega) .
$$
The Lie derivative also satisfies the relation
$$
\mathcal{L}_{f X} \omega=f \mathcal{L}_X \omega+d f \wedge i_X \omega
$$


\hypertarget{push-forward}{%
\subsubsection{Group-Valued Functions}\label{sec:push-forward}}
The \emph{push-forward} $g_*$ by a map $g: \mathcal{U} \rightarrow \mathcal{V}$ between manifolds $\mathcal{U}$ and $\mathcal{V}$ 
at a point $\boldsymbol{x}\in\mathcal{U}$ is a linear map:
\(g_*: T_{\boldsymbol{x}} \mathcal{U} \rightarrow T_{g(\boldsymbol{x})} \mathcal{V}\)
from the tangent space \(T_{\boldsymbol{x}} \mathcal{U}\) of \(\mathcal{U}\) at \(\boldsymbol{x}\)
to the tangent space \(T_{g(\boldsymbol{x})} \mathcal{V}\)
of \(\mathcal{V}\) at \(g(\boldsymbol{x})\) that is defined by:
%
\begin{equation*}
%g_*(x)\left[\dot{\gamma}(0)\right]=(g \circ \gamma)^{\cdot}(0)
g_*\left[\dot{\gamma}(0)\right]=\left.\frac{\mathrm{d}}{\mathrm{d}\varepsilon} g \circ \gamma(\varepsilon) \right|_{\varepsilon=0}
\end{equation*}
%
where \(\gamma\) is a curve in \(\mathcal{U}\) with \(\gamma(0)=\boldsymbol{x}\) and
\(\dot{\gamma}(0)\) is the tangent vector to the curve \(\gamma\) at \(0\) .
% In other words, the push-forward of the tangent vector to the curve \(\gamma\) at
% a point \(\varepsilon\)
% is the tangent vector to the curve \(g \circ \gamma\) at \(\varepsilon\).

\hypertarget{remarks02}{%
\subsection{Remarks}\label{remarks02}}
\begin{itemize}
\item
  \(\boldsymbol{s}\) and \(\boldsymbol{e}\) are regarded as functions within a single vector
  space. In subsequent developments, their derivatives only arise in
  contexts where the configuration has been fixed, thus permitting the meaningful
  use of the linear directional derivative formula in \cref{eq:dlinear}.
\end{itemize}
